\scanpage{6}
P.ex. il suffit de se limiter aux extensions de type fini du corps premier $Q$ -- on trouve un foncteur conservatif de la cat. de ces corps dans celle des groupoïdes \inserted{(ou des topos)}, au sens : un morphisme de corps qui donne une équivalence de groupoïdes (ou de Topos) est un \uline{iso}.

\margin{en car $\unsure{p>0|p\ne0}$ le foncteur est conservatif si de plus on passe à la catégorie \unsure{quot$^t$|quotient} de celle des corps où une extension radicielle finie est considérée comme triviale... (le foncteur $K\unsure{-|\longrightarrow}B_K$ ou $K\unsure{-|\longrightarrow}\Pi_K$ se factorise par cette cat. quotient.)}

Quand on prend des corps quelconques, le 2-foncteur $K\unsure{-|\longrightarrow}B_K$ ou $K\longrightarrow\Pi_K$ \inserted{ou $(K,\bar K)\unsure{-|\longrightarrow}E_{\bar K/K}$} est cependant loin d'être fidèle. Ainsi, si $K$ est sép clos, $B_K$ est le ``topos ponctuel'', $\Pi_K$ le groupoïde ponctuel, $E_{\bar K/K}\simeq(1)$ -- il est clair que les morphismes entre corps sép. clos ne sont pas décrits par les morphismes entre leurs topos étales, ou groupoïdes fondamentaux ! Pour cette raison, il y a lieu d'associer à un corps $K$ un objet plus fin que $B_K$ ou $\Pi_K$,
\margin{5}
